% ---------------------------------------------------------------
% figures/static/static/tikz/windkessel.tex
% 2-element / 3-element Windkessel RC analog.
% Terminal flow Q_Gamma(t), proximal R_p, capacitor C, distal R_d, p_out.
% Paired with cartoon of a compliant terminal vessel -> peripheral resistance.
% ---------------------------------------------------------------
\begin{tikzpicture}[fig, scale=1.0,
    wire/.style={mathblue!70!black, line width=0.7pt},
    qarrow/.style={->, bloodred!70!black, line width=1.1pt},
    capplate/.style={mathblue!70!black, line width=1pt},
    res/.style={draw=mathblue!70!black, fill=white, line width=0.7pt,
                rectangle, minimum width=10mm, minimum height=4mm,
                inner sep=1pt, font=\scriptsize}]

  % LEFT: cartoon vessel + windkessel analog
  \begin{scope}[local bounding box=cart]
    % aorta-like compliant chamber
    \path[vesselWall, rounded corners=4pt] (-0.2,1.2) rectangle (3.2,-1.2);
    \path[lumen, rounded corners=3pt] (0,1.0) rectangle (3.0,-1.0);
    % wall arrows showing distensibility
    \foreach \y in {-0.6,0,0.6}{
       \draw[<->, mathblue!70!black, line width=0.4pt]
         (1.5,\y) -- ++(0,0.32);
    }
    % flow in
    \draw[qarrow] (-1.2,0) -- (-0.05,0)
      node[midway, above, labelB, font=\small] {$Q_{\Gamma}(t)$};
    % flow out
    \draw[qarrow] (3.05,0) -- (4.2,0);
    \node[font=\small, mathblue!70!black, align=center]
      at (3.7,-0.6) {peripheral \\ bed};
    \node[font=\small\bfseries, mathblue!70!black] at (1.5,1.7)
      {compliant vessel};
  \end{scope}

  % RIGHT: circuit analog
  \begin{scope}[shift={(7.6,0)}, local bounding box=ckt]
    % nodes
    \coordinate (in) at (-1.8,0);
    \coordinate (n1) at (-0.4,0);   % after R_p
    \coordinate (n2) at (1.6,0);    % main node
    \coordinate (out) at (3.0,0);
    \coordinate (gnd) at (1.6,-1.8);

    % wires
    \draw[wire] (in) -- (-1.0,0);
    \node[res] at (-0.7,0) {$R_p$};
    \draw[wire] (-0.4,0) -- (n2) -- (out);

    % capacitor branch
    \draw[wire] (n2) -- (1.6,-0.7);
    \draw[capplate] (1.2,-0.7) -- (2.0,-0.7);
    \draw[capplate] (1.2,-0.95) -- (2.0,-0.95);
    \node[font=\scriptsize, mathblue!70!black, anchor=east] at (1.15,-0.82) {$C$};
    \draw[wire] (1.6,-0.95) -- (gnd);
    % ground symbol
    \draw[wire] (1.3,-1.8) -- (1.9,-1.8);
    \draw[wire] (1.4,-1.95) -- (1.8,-1.95);
    \draw[wire] (1.5,-2.10) -- (1.7,-2.10);

    % distal resistance branch (3-element WK): place R_d in series with C-> ground? Common 3-el: R_p in series, then (R_d || C).
    % We'll show R_d to ground (distal resistance to venous return).
    \draw[wire] (out) -- (3.6,0);
    \node[res] at (3.9,0) {$R_d$};
    \draw[wire] (4.2,0) -- (4.7,0) -- (4.7,-1.8) -- (gnd);

    % source
    \draw[qarrow] (-2.7,0) -- (in)
      node[midway, above, labelB, font=\small] {$Q_{\Gamma}(t)$};
    % p label
    \node[mathblue!70!black, font=\small\bfseries] at (n2) [above=4pt] {$p(t)$};
    \fill[mathblue!70!black] (n2) circle (1.4pt);

    % p_out reference
    \node[font=\scriptsize, mathblue] at (4.7,-2.05) [right=2pt]
        {$p_{\mathrm{out}}\!\equiv\!0$ (venous)};

    % governing ODE
    \node[font=\footnotesize, draw=mathblue!60, fill=mathblueLite,
          rounded corners=2pt, inner sep=4pt, align=left]
      at (1.3,2.1)
      {$C\dfrac{dp}{dt}+\dfrac{p-p_{\mathrm{out}}}{R_d}
         = Q_{\Gamma}(t)$\\[-1pt]
       $p_{\Gamma,\mathrm{g}}(t)=p(t)+R_pQ_{\Gamma}(t)$};
  \end{scope}

  % connecting analogy arrow
  \draw[->, mathblue!70!black, line width=0.8pt, dashed]
    (cart.east) -- (ckt.west)
    node[midway, above, labelM, font=\scriptsize, align=center]
       {0D \\ analog};
\end{tikzpicture}
